% Výsledná formule převodu NzMna na SAT. Každá část formule je podbarvena
% vlastní barvou a doplněna bublinou s vysvětlivkou (styly formulapart,
% formulaconn a formulacallout z assets/macros.tex).
\begin{tikzpicture}

    % ------------------------------------------------------------------
    % Jednotlivé části formule
    % ------------------------------------------------------------------
    \node[formulapart=blue] (nezavislost) at (0.2,0)
        {$\displaystyle\bigwedge_{ij\in E}(\neg v_i\lor\neg v_j)$};

    \node[formulapart=red] (radkyN) at (0.2,-1.5)
        {$\displaystyle\bigwedge_{i=1}^{k}\;\bigwedge_{j<j^{\prime}}(\neg x_{i,j}\lor\neg x_{i,j^{\prime}})$};

    \node[formulapart=green] (radkyA) at (0.2,-3)
        {$\displaystyle\bigwedge_{i=1}^{k}(x_{i,1}\lor x_{i,2}\lor\dots\lor x_{i,n})$};

    \node[formulapart=orange] (sloupce) at (0.2,-4.5)
        {$\displaystyle\bigwedge_{j=1}^{n}\;\bigwedge_{i<i^{\prime}}(\neg x_{i,j}\lor\neg x_{i^{\prime},j})$};

    \node[formulapart=purple] (vazba) at (0.2,-6)
        {$\displaystyle\bigwedge_{i=1}^{k}\;\bigwedge_{j=1}^{n}(\neg x_{i,j}\lor v_j)$};

    % ------------------------------------------------------------------
    % Levý sloupec se spojkami
    % ------------------------------------------------------------------
    \node[formulaconn] at (0.1,0)    {$\varphi\;=$};
    \node[formulaconn] at (0.1,-1.5) {$\land$};
    \node[formulaconn] at (0.1,-3)   {$\land$};
    \node[formulaconn] at (0.1,-4.5) {$\land$};
    \node[formulaconn] at (0.1,-6)   {$\land$};

    % ------------------------------------------------------------------
    % Bubliny s vysvětlivkami (zarovnané k nejširší části formule)
    % ------------------------------------------------------------------
    \node[formulacallout=blue, callout absolute pointer={(nezavislost.east)}]
        at (radkyA.east |- nezavislost) {Vybrané vrcholy tvoří nezávislou množinu};

    \node[formulacallout=red, callout absolute pointer={(radkyN.east)}]
        at (radkyA.east |- radkyN) {Nejvýše jedna jednička v~každém řádku};

    \node[formulacallout=green, callout absolute pointer={(radkyA.east)}]
        at (radkyA.east |- radkyA) {Alespoň jedna jednička v~každém řádku};

    \node[formulacallout=orange, callout absolute pointer={(sloupce.east)}]
        at (radkyA.east |- sloupce) {Nejvýše jedna jednička v~každém sloupci};

    \node[formulacallout=purple, callout absolute pointer={(vazba.east)}]
        at (radkyA.east |- vazba) {Očíslované vrcholy jsou vybrané};

\end{tikzpicture}
